\section{Discussion and Next Steps}
\paragraph{Contact evidence has a small meaning.}
A valid report means that a configured controller signed samples satisfying a
known fixture profile for this interaction, under its sensor-path assumptions.
It cannot tell whether a person, robot, prerecorded trace or malicious firmware
produced them. No amount of relabeling makes this equivalent to consent. A
contract requiring an unsupported guarantee must deny the action, even when
every reported contact matches.

\paragraph{Stronger primitives are complementary.}
Protected sensor firmware, attestation, authenticated acquisition, calibrated
hardware and a user-approved transaction can justify additional claims. Those
claims need their own evidence and failure analysis. They should not be inferred
from a longer contact pattern alone. A deployment may not need tactile evidence
at all when a trusted API and appropriate consent flow already authorize the
exact action.

\paragraph{What remains research.}
The shared core checks one modality per receipt. It is not yet a composition
calculus, dependency invalidation engine or selective acquisition planner. OpenAVT
remains an independent voice implementation. OpenAST's fresh display response and
OpenATT's contact report cannot be substituted merely because both are signed.
Future handoffs must preserve resource, purpose, episode, assumptions and
supporting dependencies. A useful scientific result would establish that rule
or reduce acquisition/intervention cost without weakening the strongest policy
control. The present implementation demonstrates neither advantage.

The practical next step is a bounded hardware fixture with authenticated capture
ordering, protected signing keys and a cooperating backend. Human studies and
hazardous actuation are unnecessary for the current negative-result reference;
any later consent or usability claims need an appropriate study design.

